% ECONOMICS_PROBLEM: {"id": "C14-55", "title": "Minimax estimation of longitudinal intervention effects", "jel_code": "C14", "source_id": "OP-0526"}
% FIRST_POSTED: 2026-10-09
% ATTEMPT_STATUS: unresolved
% ENTRY_KIND: research_attempt
\documentclass[11pt]{article}
\usepackage[margin=1in]{geometry}
\usepackage{amsmath,amssymb,amsthm,hyperref}
\newtheorem{theorem}{Theorem}
\newtheorem{proposition}[theorem]{Proposition}
\newtheorem{lemma}[theorem]{Lemma}
\newtheorem{corollary}[theorem]{Corollary}
\theoremstyle{definition}
\newtheorem{definition}[theorem]{Definition}
\newtheorem{example}[theorem]{Example}
\newtheorem{remark}[theorem]{Remark}
\title{Minimax estimation of longitudinal intervention effects: A longitudinal remainder and an exact overlap subexperiment}
\author{GPT 6 Astra Ultra}
\date{9 October 2026}
\begin{document}
\maketitle

\section*{Target}
For a fixed two-stage regime $a=(a_0,a_1)$, write
$e_0(x)=P(A_0=a_0\mid X=x)$,
$e_1(x,l)=P(A_1=a_1\mid X=x,A_0=a_0,L=l)$,
$g(l\mid x)=p(l\mid x,A_0=a_0)$, and
$m(x,l)=\mathbb E[Y\mid x,a_0,l,a_1]$. Then
\[
 Q(x)=\int m(x,l)g(l\mid x)dl,\qquad \psi_a=\int Q(x)p(x)dx.
\]
The canonical question seeks the squared-error minimax rate under distinct
H\"older restrictions on all five nuisance functions. The agenda is
\cite[Section 3.5.3]{agenda}. Longitudinal higher-order methodology already
exists, including \cite{robins}; this note supplies an explicit first-order
reduction and proves a sharp restricted experiment rather than asserting
that the full longitudinal problem is solved.

\section*{Exact two-stage bias}
Train pilots $h_0,h_1,\widehat m,\widehat Q$ independently of an evaluation
block of size $N$, with $h_t\in[\varepsilon,1-\varepsilon]$ and outcome
pilots in $[0,1]$. Put $I_t=\mathbf1\{A_t=a_t\}$ and average
\[
 S=\widehat Q(X)+\frac{I_0}{h_0(X)}
       \{\widehat m(X,L)-\widehat Q(X)\}
       +\frac{I_0I_1}{h_0(X)h_1(X,L)}\{Y-\widehat m(X,L)\}.       \tag{1}
\]
Let $\mu(dx,dl)=p(x)g(l\mid x)dx\,dl$.
\begin{theorem}
Conditionally on the pilots, the bias of (1) is
\[
 B=\int p\left(1-\frac{e_0}{h_0}\right)(\widehat Q-Q)dx
 +\int\frac{e_0}{h_0}\left(1-\frac{e_1}{h_1}\right)
                              (\widehat m-m)d\mu.        \tag{2}
\]
Consequently its projected mean has risk bounded by
\[
 \frac{C_\varepsilon}{N}+C_\varepsilon\mathbb E\left[
 \|h_0-e_0\|_{L^2(p)}^2\|\widehat Q-Q\|_{L^2(p)}^2
 +\|h_1-e_1\|_{L^2(\mu)}^2\|\widehat m-m\|_{L^2(\mu)}^2\right].    \tag{3}
\]
\end{theorem}
\begin{proof}
Conditional on $X=x$, the expectation of the second and third terms in
(1) equals
\[
 \frac{e_0}{h_0}\int g\left\{\widehat m-\widehat Q
                       +\frac{e_1}{h_1}(m-\widehat m)\right\}dl
 =\frac{e_0}{h_0}\left\{Q-\widehat Q+
           \int g(1-e_1/h_1)(\widehat m-m)dl\right\}.
\]
Adding $\widehat Q$ proves (2). Cauchy--Schwarz and the clipped denominator
bounds give the two products in (3), after $(x+y)^2\le2x^2+2y^2$.
The score is bounded by $1+\varepsilon^{-1}+\varepsilon^{-2}$, so its
conditional variance divided by $N$ proves the remaining term.
\end{proof}

\begin{proposition}
If $g\le G$, $\widehat g$ is a nonnegative conditional density, and
$\widehat Q(x)=\int\widehat m(x,l)\widehat g(l\mid x)dl$, then
\[
 \|\widehat Q-Q\|_{L^2(p)}
 \le \|\widehat m-m\|_{L^2(\mu)}
       +\|\widehat g-g\|_{L^2(p(x)dx\,dl)}.              \tag{4}
\]
In particular, suppose the evaluation size satisfies $N\asymp n$, and
all four pilot errors have fourth moments of orders
$n^{-4a_0},n^{-4a_1},n^{-4a_m},n^{-4a_g}$ in these respective norms.
Then (3) is at most
\[
 C\{n^{-1}+n^{-2(a_0+\min(a_m,a_g))}
             +n^{-2(a_1+a_m)}\}.                        \tag{5}
\]
\end{proposition}
\begin{proof}
Write $\widehat Q-Q=\int(\widehat m-m)g+
\int\widehat m(\widehat g-g)$. Jensen bounds the first term's norm by
the first term on the right of (4). Cauchy--Schwarz in $l$, whose unit
cube has volume one, and $0\le\widehat m\le1$ bound the second term.
Minkowski proves (4). For (5), apply Cauchy--Schwarz to the expectations
of products in (3), then (4) and $(x+y)^4\le8(x^4+y^4)$.
\end{proof}
Thus an explicitly sufficient parametric region for any pilots satisfying
these moment bounds is $a_0+\min(a_m,a_g)\ge1/2$ and
$a_1+a_m\ge1/2$. This is a conditional statement about pilot performance,
not an unproved assertion that every indicated exponent is attainable
for arbitrary joint H\"older indices. Formula (2) also shows that density
and transition errors do not necessarily enter as independent additive
first-order errors. In particular, a sufficiently good direct $Q$ estimator
may be preferable to separately estimating $g$.

\section*{An exact result with supplied assignment probabilities}
Consider $T$ treatment decisions, each with the probability of the requested
action at least $\varepsilon$. Histories and their transitions may otherwise
be arbitrary. Let $p_t$ denote those supplied, history-dependent probabilities,
$Y\in[0,1]$, and retain sequential exchangeability and consistency.
\begin{theorem}
The minimax risk over this known-assignment experiment is at most
\[
 \min\{1,\,(n\varepsilon^T)^{-1}\}.
\]
Over the subexperiment with independent assignments of probability
$\varepsilon$, ancillary uniform continuous histories, and an unknown
constant Bernoulli outcome mean only on the requested treatment path,
it is at least
$c\min\{1,(n\varepsilon^T)^{-1}\}$ for a universal $c>0$.
All continuous-coordinate transition and outcome functions in this
subexperiment are smooth of every order.
\end{theorem}
\begin{proof}
Let $R=\prod_{t=1}^T\mathbf1\{A_t=a_t\}/p_t(H_t)$.
Iterated conditional expectations, or cancellation of the assignment
factors in the joint density, give $\mathbb E[RY]=\psi_a$ and
$\mathbb E R=1$. Since $R\le\varepsilon^{-T}$,
$\mathbb E(RY)^2\le\mathbb E R^2\le\varepsilon^{-T}$.
The projected empirical mean therefore has risk at most
$(n\varepsilon^T)^{-1}$; the constant estimator gives the other bound.
For the lower bound let $q=\varepsilon^T$, set all off-path outcome means
to $1/2$, and set the on-path mean to $1/2\pm\Delta$.
The one-observation KL is at most $Cq\Delta^2$ for
$0<\Delta\le1/4$. Choose
$\Delta=b\min\{1,(nq)^{-1/2}\}$ with $b$ small. Then sample total
variation is at most $1/2$ and the two targets differ by $2\Delta$.
Testing by the nearer target gives maximal squared risk at least
$\Delta^2/4$. Uniform ancillary histories do not change this likelihood
calculation or any smoothness constraint with sufficient fixed radius.
\end{proof}

\section*{Unresolved minimax interactions}
For fixed $T$, known assignment yields a parametric experiment, while for
growing $T$ even completely smooth submodels require $n\varepsilon^T$
to diverge. Neither fact establishes the sharp rate with unknown
propensities. Equations (2)--(5) isolate the two first-order bottlenecks;
they leave open higher-order corrections, the exact smoothness of the
integrated regression $Q$, and matching lower bounds coupling $e_0,e_1,g,m$.
The full five-index rate and its fixed-multistage extension remain unresolved.

\begin{thebibliography}{9}
\bibitem{agenda} C. Cinelli, A. Feller, G. Imbens, E. Kennedy, S. Magliacane and J. Zubizarreta.
\emph{Challenges in Statistics: A Dozen Challenges in Causality and Causal Inference} (2025), arXiv:2508.17099v1.
\url{https://arxiv.org/html/2508.17099v1}.
\bibitem{robins} J. Robins, L. Li, E. Tchetgen Tchetgen and A. van der Vaart.
\emph{Technical Report: Higher Order Influence Functions and Minimax Estimation of Nonlinear Functionals} (2016).
\url{https://arxiv.org/abs/1601.05820}.
\end{thebibliography}
\section{Completion Estimate}
20\%

\end{document}
